\documentclass[10pt]{article}

% ---------- PACKAGES ----------
\usepackage[margin=1.5cm]{geometry}
\usepackage{amssymb, amsthm}
\usepackage{mathtools}
\usepackage{fancyhdr}
\usepackage{enumitem}
\usepackage{xcolor}
\usepackage{setspace}
\usepackage{tcolorbox}
\usepackage[hidelinks]{hyperref}

% ---------- CUSTOM COUNTERS & COMMANDS ----------
\newcounter{problemnum}
\newcounter{subpartnum}[problemnum]

\newcommand{\problem}[1]{
  \stepcounter{problemnum}
  \setcounter{subpartnum}{0}
  \vspace{0.8em}
  \noindent
  {\boldmath{\textbf{#1}}}
  \nopagebreak
  \vspace{0.4em}
  \par
}

\newcommand{\problemtag}[2]{%
    \noindent
    \colorbox{black}{%
        \textcolor{white}{\textbf{#1 (pg.\ #2)}}%
    }
    \addcontentsline{toc}{subsubsection}{Exer. #1}%
}

\newcommand{\subpart}{
  \stepcounter{subpartnum}
  \vspace{0.3em}
  \noindent
  \textbf{(\alph{subpartnum})}
  \quad
}

\newenvironment{solution}{
  \vspace{0.2em}
  \noindent
  \normalsize
  \textit{Solution:}
  \nopagebreak
  \begin{list}{}{
    \setlength{\leftmargin}{2em}
    \setlength{\parsep}{0.8em}
    \setlength{\itemindent}{0em}
  }
  \item
}{
  \end{list}
  \vspace{0.3em}
}

% BOLD MATH STATEMENT COMMAND
\newcommand{\boldstatement}[1]{
  {\boldmath \underline{\textbf{#1}}}
}

% RIGHT-SIDE JUSTIFICATION COMMAND
\newcommand{\justify}[1]{
  \hfill \textit{(#1)}
}

% Alternative: For multi-line proofs with justifications
\newcommand{\justifyline}[2]{
  #1 \hfill \textit{(#2)} \\
}

% ---------- GLOBAL STYLE ----------
\setlength{\parindent}{0pt}
\setlength{\parskip}{0.5em}
\pagestyle{fancy}
\fancyhf{}
\cfoot{\thepage}
\renewcommand{\headrulewidth}{0pt}

% ---------- DOCUMENT INFO ----------
\newcommand{\courseName}{MATB24H3}
\newcommand{\assignmentNum}{1}
\newcommand{\studentName}{Your Name}
\newcommand{\studentNumber}{Your Student Number}
\newcommand{\dueDate}{Due Date}

% ---------- TITLE ----------
\title{\textbf{\courseName}}
\author{\studentName}

% ----- SHORTHAND ----------
\newcommand{\Q}{\mathbb{Q}}
\newcommand{\Z}{\mathbb{Z}}
\newcommand{\R}{\mathbb{R}}
\newcommand{\C}{\mathbb{C}}
\newcommand{\F}{\mathbb{F}}
\newcommand{\N}{\mathbb{N}}
\newcommand{\K}{\mathbb{K}}
\newcommand{\cX}{\mathcal{X}}
\newcommand{\cM}{\mathcal{M}}
\newcommand{\cR}{\mathcal{R}}
\newcommand{\cP}{\mathcal{P}}
\newcommand{\cL}{\mathcal{L}}
\newcommand{\cC}{\mathcal{C}}

\begin{document}
\onehalfspacing


\begin{center}
  \LARGE\textbf{MATB24H3}\\
  \Large\textbf{Exercise Questions}
\end{center}

\textbf{Student Name: \studentName} \hfill \textbf{Student Number: \studentNumber}
\vspace{0.5em}
\hrule

\tableofcontents
\newpage

% ---------- USAGE INSTRUCTIONS ----------
% Please be responsible and don't include any solutions in this file when uploading.
%
% You can use this file to write your own solutions in a separate file.
%
% Please use the features \problems, \subpart, \begin{solution}, \justify, \boldstatement{}, etc. to write your solutions in a clear and organized manner.
%
% Please help gather all the problems for the course into this file for the good of the public.
%
% The following are shorthand for the following commands:
%
%   \mathbb{Q} => \Q
%   \mathbb{Z} => \Z
%   \mathbb{R} => \R
%   \mathbb{C} => \C
%   \mathbb{F} => \F
%   \mathbb{N} => \N
%   \mathbb{K} => \K
%   \mathcal{M} => \cM
%   \mathcal{R} => \cR
%   \mathcal{P} => \cP
%   \mathcal{L} => \cL
%   \mathcal{C} => \cC

% ---------- PROBLEMS ---------- 
\section{Notation/Terminologies}
...

\section{Fields \& Rings \& Groups}

\subsection{Preliminary}
\problem{
  \problemtag{2.1.1.}{6} Consider the set $X = \Q \setminus \{0\}$ and the map $\star$ is depicted by $a \star b = \frac{a}{b}$. Is $\star$ associative?
}
\begin{solution}
  % ==== Place for solution here ====
\end{solution}

\problem{
  \problemtag{2.1.2.}{7} Is $(\cP(\N), \cap)$ a monoid? What is the identity element? How about $(\cP($\textit{X}$), \cup)$ for any arbitrary set $X$?
}
\begin{solution}
  % ==== Place for solution here ====
\end{solution}


\subsection{Fields}
\problem{
  \problemtag{2.2.1}{13} Suppose $d$ is square-free, so there is no integer $m$ such that $m^2 = d$. Show that \\ ($\Q(\sqrt{d})$, +, ·)
forms a field. What happens if $d = m^2$ for some $m$?
}
\begin{solution}
  % ==== Place for solution here ====
\end{solution}

\problem{
  \problemtag{2.2.2.}{12} Now, can we extend a \textit{finite field} in the same way we extended an infinite field? If so, give an example.
}
\begin{solution}
  % ==== Place for solution here ====
\end{solution}

\problem{
  \problemtag{2.2.3.}{14} Prove that $0_\F$, $1_\F$ and additive inverse element $-a$ of any element $a \in \F$ are unique.
}
\begin{solution}
  % ==== Place for solution here ====
\end{solution}


\newpage
\problem{
  \problemtag{2.2.4.}{15} Prove $(-a)(-b) = ab$ for all $a, b \in \F$.
}
\begin{solution}
  % ==== Place for solution here ====
\end{solution}


\problem{
  \problemtag{2.2.5.}{16} Prove that the triple $(\Z_p, + \mod p, \cdot \mod p)$ is a field $\iff$ p is a prime. Use the fact that if two integers $m$, $n$ are relatively prime, then there
exist integers $r$ and $s$ such that $mr + ns = 1$.
}
\begin{solution}
  % ==== Place for solution here ====
\end{solution}


\problem{
  \problemtag{2.2.6.}{20} Let, 
  \[ S = \{ r^0 := 1, r, r^2, r^3 \} \]
  Prove that the 2-tuple $(S, \circ)$, where $S$ is the set of symmetries of a square and $\circ$ is the composition operation, does indeed form a group. Moreover, find the symmetry
group of a rectangle.
}
\begin{solution}
  % ==== Place for solution here ====
\end{solution}


\problem{
  \problemtag{2.2.7.}{22} Characterize all idempotent elements of $(\cM_{2\times 2}(\R), +, \cdot)$.
}
\begin{solution}
  % ==== Place for solution here ====
\end{solution}


\problem{
  \problemtag{2.2.8.}{26} Prove that $\cC(\cR)$ is a subring of any ring $(\cR, +, \cdot)$.
}
\begin{solution}
  % ==== Place for solution here ====
\end{solution}


\problem{
  \problemtag{2.2.9.}{26} Let $(\cR, +, \cdot)$ be a ring and suppose that $\alpha \in \cR$ has a multiplicative inverse. Moreover,
suppose that $\alpha \in C(\cR)$. Is $\alpha^{-1} \in C(\cR)$?
}
\begin{solution}
  % ==== Place for solution here ====
\end{solution}


\problem{
  \problemtag{2.2.10.}{28} Does the set of functions with infinitely many roots form a subring of $\mathcal{F}(\R)$? Justify your answer.
}
\begin{solution}
  % ==== Place for solution here ====
\end{solution}

\problemtag{[Refer to notes]}{29}
\newpage
\problem{
  \problemtag{2.2.12.}{29} Determine the zero and the identity elements of the polynomial ring $(\cR[x], +, \cdot)$ where $\cR = \cC(\cM_{2\times 2}(\R))$.
}
\begin{solution}
  % ==== Place for solution here ====
\end{solution}


\problem{
  \problemtag{2.2.13.}{32} True/False. Justify:
  \begin{enumerate}
    \item Let $\F$ be any field. Suppose that $\alpha$ and $\beta$ are two elements of an extension field $\K$ of $\F$. Then $\F(\alpha)$ and
$\F(\beta)$ are always distinct extension fields of $\F$.
    \item Given any set $X$, there is no operation · that can be defined on $[X]^{\leq \left|X\right|}$ that makes it a group.
    \item Given a monoid $(\mathcal{S}, *)$ there exists a binary operation $\cdot$ on $\mathcal{S}$ such that $(\mathcal{S}, \cdot)$ no longer admits an identity element. If so, what assumption did you need to make?
    \item $(\R, \cdot)$ forms a group.
    \item $(\R, +)$ forms a group with the identity element being 1.
    \item Suppose $(\F, +, \cdot)$ is the field. Then $(\F, +)$ and $(\F, \cdot)$ both form a group.
  \end{enumerate}
}
1) \begin{solution}
  % ==== Place for solution here ====
\end{solution}
2) \begin{solution}
  % ==== Place for solution here ====
\end{solution}
3) \begin{solution}
  % ==== Place for solution here ====
\end{solution}
4) \begin{solution}
  % ==== Place for solution here ====
\end{solution}
5) \begin{solution}
  % ==== Place for solution here ====
\end{solution}
6) \begin{solution}
  % ==== Place for solution here ====
\end{solution}


\newpage

\section{Vector Spaces}
\subsection{Basis}
\problem{
\problemtag{3.1.1.}{37} Had we defined vector addition by \[\hat{A}\boxplus\hat{B}=[a_{ij}\cdot b_{ij}] \text{ i.e., pointwise multiplication of the entries}\]
in Example 3.1.3, would 
\[(\mathcal{M}_{m\times n}(\F),\,\boxplus,\,\odot,\,(\mathbb{F},+,\cdot)) \]
form a vector space over $\mathbb{F}$?
}
\begin{solution}
  % ==== Place for solution here ====
\end{solution}


\problem{
\problemtag{3.1.2.}{40} Prove properties 2, 4 and 5 of Theorem 3.1.2, namely,
\begin{tcolorbox}
For any vector space $\mathbb{V}=(V,\boxplus,\odot,(\mathbb{F},+,\cdot))$ over field $\mathbb{F}$:
\begin{itemize}
    \item 2) The additive inverse of any $v\in V$ is unique.
    \item 4) Any scalar of the zero vector yields the zero vector, i.e. $\forall\,\alpha\in\mathbb{F},\,a\odot\vec{0}=\vec{0}$.
    \item 5) For all $v\in V$ and $c\in\mathbb{F}$, $(-c) \odot v = c \odot (-v) = -1 \odot (c \odot v)$.
\end{itemize}
\end{tcolorbox}
}

\begin{solution}
  % ==== Place for solution here ====
\end{solution}


\subsection{Subspaces}
\problem{
\problemtag{3.2.1.}{42} Prove Theorem 3.2.2, namely
\begin{tcolorbox}
If $S\subseteq V$ then $\text{span}(S)\subseteq V$ is a subspace of $V$.
}
\end{tcolorbox}
\begin{solution}
  % ==== Place for solution here ====
\end{solution}


\newpage
\problem{
\problemtag{3.2.2.}{42} Prove Theorem 3.2.3, namely:
\begin{tcolorbox}
Let $J$ be any index set (It is a set whose elements are used to enumerate the elements of set). Let $\{W_\alpha\mid\alpha\in J\}$ be a collection of subspaces of $V$. Then $\bigcap\limits_{\alpha\in J} W_{\alpha}$ is a subspace of $V$.
\end{tcolorbox}
}
\begin{solution}
  % ==== Place for solution here ====
\end{solution}

\problem{
\problemtag{3.2.3.}{43} Prove Theorem 3.2.4, namely: 
\begin{tcolorbox}
Let $W_1\subseteq W_2\subseteq W_3\subseteq\dots$ be an increasing chain of subspaces of $V$. Then $\bigcup\limits_{n\in\mathbb{N}}W_n$ is a subspace.
\end{tcolorbox}
}
\begin{solution}
  % ==== Place for solution here ====
\end{solution}


\subsection{Bases}
\problem{
\problemtag{3.3.1.}{46} Show that the set $\{e^x,\,e^{2x},\,e^{3x}\}$ is linearly independent in the vector space of $\mathcal{F}(\mathbb{R})$ over $\mathbb{R}$.}
\begin{solution}
  % ==== Place for solution here ====
\end{solution}


\newpage
\subsection{Linear Transformations}
\problem{
\problemtag{3.4.1.}{50} Prove Proposition 3.4.1, namely, 
\begin{tcolorbox}
The following: Let $T\colon V\to W$ be a linear transformation from $V$ to $W$. Then:
\begin{enumerate}
    \item For all $x,\,y\in V$, $T(x-y)=T(x)-T(y)$
    \item T maps the zero vector $\vec{0}_V$ of $V$ to the zero vector $\vec{0}_W$ of W.
    \item If $T'\colon W\to Z$ is also a linear transformation, their composition $T'\circ T$ is also a linear transformation from $V$ to $Z$ where $T'\circ T=T'(T(v))$.
    \item If $V_1$ is a subspace of $V$ and $W_1$ is a subspace of $W$, then $T[V_1]$ is a subspace of $W$ and $T^{-1}[W_1]$ is a subspace of $V$.
\end{enumerate}
\end{tcolorbox}
}
1) \begin{solution}
  % ==== Place for solution here ====
\end{solution}
2) \begin{solution}
  % ==== Place for solution here ====
\end{solution}
3) \begin{solution}
  % ==== Place for solution here ====
\end{solution}
4) \begin{solution}
  % ==== Place for solution here ====
\end{solution}


\problem{
\problemtag{3.4.2.}{51}Prove Theorem 2.4.2 
\begin{tcolorbox}
(currently, this theorem does not exist on the lecture notes, but Pourya will surely fix this in no time.)
\end{tcolorbox}
}
\vspace{0.8em}
\begin{solution}
  % ==== Place for solution here ====
\end{solution}


\newpage
\problem{
\problemtag{3.4.3.}{52} Prove Theorem 3.4.2, namely, 
\begin{tcolorbox}
show that $(\cL(V,W),\boxplus,\odot,(\F,+\cdot))$ forms a vector space over $\F$, where $\boxplus$ and $\odot$ are the pointwise addition of linear maps and scalar multiplication with the linear map in $W$, respectively. That being, $\boxplus$ is the vector addition operation in $W$ and $\odot$ is the scalar multiplication operation in $W$.
\end{tcolorbox}
}
\begin{solution}
  % ==== Place for solution here ====
\end{solution}


\problem{
\problemtag{3.4.4.}{53} Prove that a linear map $T\colon V\to W\text{ is invertible }\iff T\text{ is injective and surjective}$.}
\begin{solution}
  % ==== Place for solution here ====
\end{solution}

\problem{
\problemtag{3.4.5.}{53}Let $V$ and $W$ be two vector spaces over a field $\F$ with dimensions $n$ and $m>3$ respectively. Construct three different linearly independent linear maps from $V$ to $W$.}
\begin{solution}
  % ==== Place for solution here ====
\end{solution}


\problem{
\problemtag{3.4.6.}{53} Prove whether or not the triple $(\cL(\R,\,\R),\,+,\,\cdot)$ forms a subring of $\mathcal{F}(\R)$ under point-wise addition and multiplication.}
\begin{solution}
  % ==== Place for solution here ====
\end{solution}



% ========== START HERE ==========
\end{document}